\chapter{Introduction : Méthodologie de correspondance}

Cette troisième partie établit les correspondances explicites entre mon parcours professionnel, les expériences détaillées en Partie 2, et le référentiel RNCP38152 du Master MEEF Pratiques et Ingénierie de la Formation (PIF) \parencite{rncp38152}. Elle démontre que mes acquis couvrent l'intégralité des 9 blocs de compétences attendus pour l'obtention du diplôme par validation des acquis de l'expérience \parencite{vademecum_vae_inspe}.

\section{Cadre de référence : le RNCP38152}

Le référentiel RNCP38152 « MASTER - Métiers de l'enseignement, de l'éducation et de la formation (MEEF), pratiques et ingénierie de la formation » définit les compétences attendues d'un expert en ingénierie de formation à un niveau de qualification 7 (Bac+5) \parencite{rncp38152,referentiel_meef}.

Ces compétences s'organisent en 9 blocs couvrant l'ensemble des dimensions de l'ingénierie de formation, de la conception à l'évaluation, en passant par le pilotage de projets et l'analyse réflexive \parencite{competences_transversales}.

\section{Méthodologie d'analyse}

Pour chaque bloc de compétences, je présenterai :

\begin{enumerate}[leftmargin=*]
    \item \textbf{La définition officielle} du bloc selon le référentiel RNCP38152
    \item \textbf{Les compétences détaillées} attendues
    \item \textbf{Mes expériences correspondantes} : références précises aux situations décrites dans la Partie 2 et mon parcours (Partie 1)
    \item \textbf{Les preuves tangibles} : productions, réalisations, attestations
    \item \textbf{Le niveau de maîtrise atteint} : indication du degré d'expertise développé
\end{enumerate}

Cette approche systématique permet au jury VAE d'évaluer objectivement l'adéquation entre mes acquis et les attendus du diplôme \parencite{vademecum_vae_inspe}.
